\subsubsection{Alternating sequences with larger jumps and the finite-class length}
\label{sec:inert-thick-word}
\hypertarget{source-inert-thick-word}{}

For an alternating sequence whose larger jumps dominate, the finite-class estimate needs information about the image of each successive lattice map. We follow a lifted class through the sequence and control the restrictions on its coefficients at the even positions.

\begin{established}[Marked alternating word]\label{hyp:alternating-thick-word}
Let \(F/\Q_\ell\) be unramified of even degree \(n\), with
\(\ell\ge11\). Let \(f:T_1\to T_2\) be the marked crystalline surjection of
free stable lattices over the coefficient ring \(O\).
Take its field \(K_O\) to contain the embeddings of \(F\) and
the primitive \(\ell\)-th roots. The descent
of the finite-class bound to the original coefficients uses the same
faithfully flat base change as in Section~\ref{sec:inert-long-orbit}.
Write \(r=\rk_O T_2\). Index the embeddings cyclically and absorb the
parity shift in the choice of origin. If \(J_i\) denotes the source jumps at letter \(i\), their boxes are
\[
 J_i\subseteq[-\eta,\eta]\cup\begin{cases}[-b_i-S,-b_i],&i\text{ even},\\ {}[b_i,b_i+S],&i\text{ odd}.\end{cases}
\]
At most \(m\) target jumps, counted
with multiplicity, lie in the central interval at each embedding.
The target jumps are among the source jumps by strictness of the
crystalline filtration.

Put
\[
\begin{gathered}
 h=\max_i b_i+S,\quad H_{\rm loc}=h+\eta,\quad
 L_0=1+\log_\ell(2H_{\rm loc}),\\
 A_0=\left\lceil\frac{(2\ell-1)H_{\rm loc}}{(\ell-1)^2}
                +3+L_0\right\rceil .
\end{gathered}
\]
In the retained alternating case one has
\(b_i\ge h/2\) and
\(b_i\ge 2(A_0^{\rm real})+\eta+1\), where \(A_0^{\rm real}\) is
the expression inside the ceiling. Define
\[
 c_D=1+\frac{H_{\rm loc}+1}{\ell-1},\qquad
 C_D=rS+2mc_D .
\]
For an \(O\)-module of finite length, we use its length divided by
\(k_O=\operatorname{length}_O(O/\ell O)\).
\end{established}
\begin{proof}
At an inert prime the two embeddings of the distinguished quadratic field are exchanged by Frobenius. On an orbit carrying larger jumps at every position, Lemma~\ref{lem:marked-signed-jumps} therefore gives the displayed signs after a cyclic reindexing. The central multiplicity is the number of central columns in the marked rank-\(r\) weights. The target inherits the source support by strict exactness. On the retained word
\(b_0\le b_i\le b_0+\eta\) with \(b_0=b+O(t)\), whereas
\(S+\eta+\log H=o(b_0)\) uniformly in the host. Thus
\(h\le b_0+o(b_0)\), and the coefficient
\(2(2\ell-1)/(\ell-1)^2\le 42/100\) for \(\ell\ge11\)
gives the second inequality beyond a common threshold.
In particular \(b_i>\eta\).
\end{proof}

\paragraph{Passing from a chain to a length.}
Set \(M_2=H^1_f(F,T_2)\) and \(I=f_*H^1_f(F,T_1)\).
Choose a free \(O\)-module \(A\) of rank at most \((n+1)r\)
with a surjection \(A\to M_2\), using the generator argument in
Section~\ref{sec:inert-long-orbit}. Representatives for a basis
give a linear map to the finite cocycles themselves.
\begin{lemma}[Length from the successive restrictions]
Let \(\ell^{e'}M_2\subset I\). Suppose the final submodule of a
chain in \(A\) maps into \(I\) after multiplication by \(\ell^e\).
Let \(N_g,N_d\) be the numbers of the generator-bounded and length-bounded successive quotients, respectively. If each of the former requires at most \(m\) generators and each of the latter has normalized length at most \(C_D\), the numerical hypotheses
\[
 N_g\le4n,\qquad N_d\le3n/2
\]
give the bound \eqref{eq:thick-chain-mass} on \(\operatorname{len}(M_2/I)\).
\end{lemma}
\begin{proof}
Suppose \(A=A^{(0)}\supset\cdots\supset A^{(k)}\) and
\(\ell^e A^{(k)}\) maps into \(I\).
Suppose also that \(\ell^{e'}M_2\subset I\).
If a successive quotient of the chain is generated by at most
\(m\) elements, its image in the filtration of \(M_2/I\) costs at
most \(me'\) in normalized length. If a successive quotient has
normalized length at most \(C_D\), the corresponding image costs
at most \(C_D\). Indeed these successive images are quotients of
the successive quotients of the chain.

At the last term the image
is killed by \(\ell^e\) and has at most \(\rk_O A\) generators,
since \(A^{(k)}\) is free of rank at most \(\rk_O A\).
Thus \(4n\) steps of the first kind and \(3n/2\) of the second
give normalized length at most
\begin{equation}\label{eq:thick-chain-mass}
           (n+1)r e+4nm e'+\tfrac32 n C_D .
\end{equation}
For later use, a different sufficient condition is
\(\ell^e M_2\subset I+\langle Z\rangle_O\), where \(Z\) is a finite
set of classes in \(M_2\). The normalized length is then at most
\((n+1)r e+|Z|e'\). Take the quotient by the image of
\(\langle Z\rangle_O\) first, and use \(\ell^{e'}M_2\subset I\)
on that image. These arguments apply to the full integral finite
group, including its torsion.

\end{proof}
\paragraph{Wach coordinates and the initial finite classes.}
\label{par:thick-initial-classes}
Use the tame-invariant parameter \(X\) and \(q=X+\ell\) of
Paragraph~\ref{par:wach-cut-correction}, and put \(p_j=\phi^j(q)\).
Write \(N_i(T)\) for the tame-invariant lattice, \(N(T)\) for their
product, and \(P_i^{(j)}(T)=\langle\Phi^jN_i(T)\rangle\) for its
integral Frobenius span at \(i+j\). Let \(a=A_0\),
\(B_i=b_i+S\) and \(H=H_{\rm loc}\). The two signed height bounds
at a destination \(i\) are
\[
 q^{M_i}P_{i-1}^{(1)}(T)\subset N_i(T),\qquad
 q^{H_i}N_i(T)\subset P_{i-1}^{(1)}(T),\qquad
 (M_i,H_i)=
 \begin{cases}(B_i,\eta)&i\text{ even},\\
 (\eta,B_i)&i\text{ odd}.\end{cases}
\]
Use the cyclotomic lattice comparison over the unramified field
and the filtered comparisons at the primitive roots
\cite[Proposition III.2.1 and Corollary III.4.5]{berger-wach}.
\begin{lemma}[Annihilators at the initial letters]
For the given surjection, and for the crystalline extension of
the trivial lattice by the target with its quotient map, let \(C_i\) be the cokernel on the \(i\)-th tame-invariant letter. Then
\[
 q^a C_i=\ell^R C_i=0,\qquad R=\left\lfloor\frac{\ell H}{(\ell-1)^2}+3+\log_\ell(H+1)\right\rfloor.
\]
\end{lemma}
\begin{proof}
The following image bound uses the maximum \emph{letter width}
\(H=\max_i(M_i+H_i)\). Do the calculation before tame invariants
in the ordinary parameter \(\pi\).
The common character with labelled filtration jump \(+M_i\) shifts both lattices by
\[
 [-M_i,H_i]+M_i=[0,M_i+H_i]\subseteq[0,H].
\]

The jumps of the decreasing filtration are now in \([0,H]\),
as required in Berger's positive comparison. Such a character is a product
of powers of the labeled Lubin--Tate characters of the
unramified field, with values in the coefficient field.
\cite[Proposition 2.3]{berger-lubin-tate}
Tensoring by this rank-one lattice preserves the annihilator
of the Wach cokernel. Rational Wach comparison is tensor
compatible. Integrally, multiplication by a generator of the
rank-one Wach lattice identifies the integral \'etale modules,
since the generator is a basis after \'etale scalar extension.
Taking their intersections with the rational Wach modules gives
the tensor identification on integral Wach lattices.
\cite[Proposition III.4.2]{berger-wach}

Apply Berger's positive comparison and the faithfully flat
descent in Paragraph~\ref{par:wach-cut-colength} after the common
twist. For the untwisted Wach cokernel \(C\), the common rank-one twist preserves the annihilator and hence gives
\[
 \pi^H C=0.
\]

The argument of that paragraph which lowers the \(\pi\)-degree
by the divided cyclotomic action gives a scalar exponent at most
\(\lfloor\ell H/(\ell-1)^2\rfloor\). The two input lattices also
occur as the extension of the trivial lattice by the target
and its trivial quotient. The jumps of the extension are those
of the target together with zero, so the same boxes apply.

After tame averaging the annihilators on each letter include
\(X^{\lfloor H/(\ell-1)\rfloor+1}\) and \(\ell^R\), with the following integer exponent
\[
 R=\left\lfloor \frac{\ell H}{(\ell-1)^2}
                 +3+\log_\ell(H+1)\right\rfloor .
\]
Set \(d=\lfloor H/(\ell-1)\rfloor\), so the two annihilators are
\(X^{d+1}\) and \(\ell^R\). In the expansion of \((X+\ell)^a\),
each term with \(X\)-degree \(k\le d\) has \(\ell\)-exponent
\(a-k\ge R\); every other term is divisible by \(X^{d+1}\).
Thus \(q^a\) annihilates the same cokernel.

The bounds are uniform in the number of letters,
the rank and the ramification of the coefficient ring.
\label{cor:inert-thick-initial-1}

\end{proof}

\begin{lemma}[Integral Herr representatives for the basis extensions]
For every basis cocycle \(y\) the two annihilators give vectors \(\mu_i(y),\nu_i(y)\) satisfying \eqref{eq:thick-extension-data} on each letter. On the free span, the choice is extended by
\[
 (\mu_i,\nu_i)\!\left(\sum_j b_j y_j\right)=\sum_j b_j\big(\mu_i(y_j),\nu_i(y_j)\big).
\]
\end{lemma}
\begin{proof}
For each basis cocycle in the free family take its extension
\(0\to T_2\to Y\to O\to0\). Its rational extension is crystalline
by the definition of the finite group. Use exactness of the
integral \((\varphi,\Gamma)\)-module functor and the Wach
intersection with the kernel. Tame averaging, of order
\(\ell-1\), preserves integral lifts.

Choose a tame-invariant
lift \(E\) of the trivial generator in the \'etale module of
\(Y\), and take the Herr pair \(y=(\Phi-1)E\),
\(z=(\gamma-1)E\) in that of \(T_2\). Denote their letter
components by \(y_i,z_i\), and use the convention
\((\gamma-1)y=(\Phi-1)z\).
The two annihilators provide Wach lifts \(E_i,E_i'\) of
\(\ell^R\) and \(q^a\) times the trivial generator in letter \(i\).
Set
\[
 \mu_i=E_i-\ell^RE|_i,\qquad \nu_i=E_i'-q^a E|_i,\qquad
 A_i=\ell^R y_i+\Phi\mu_{i-1}-\mu_i .
\]
Both differences defining \(\mu_i,\nu_i\) belong to the target
\'etale module. The Wach intersection and the first signed height
bound give
\begin{equation}\label{eq:thick-extension-data}
\begin{gathered}
 \ell^R z_i+(\gamma-1)\mu_i\in XN_i(T_2),\qquad
 q^{M_i} A_i\in N_i(T_2),\\
 \ell^R q^{H_i} A_i\in P_{i-1}^{(1)}(T_2),\qquad
 \ell^R\nu_i-q^a\mu_i\in N_i(T_2).
\end{gathered}
\end{equation}

For the third inclusion, \(A_i=\Phi E_{i-1}-E_i\) in the
extension, so \(q^{H_i}A_i\) belongs to its Frobenius span.
Expand in a free basis of \(O[[X]]\) over \(\phi(O[[X]])\),
as in the proof of the arc estimate. Each quotient coefficient
of the resulting Wach vectors vanishes by independence.
Thus the Frobenius span meets the target in the target span.
This even gives the inclusion before multiplication by \(\ell^R\).
The data in \eqref{eq:thick-extension-data} now extend linearly in the free family,
even though the individual extensions \(Y\) were chosen only
on its basis.

\end{proof}
\paragraph{First restriction at an even letter.}
\begin{lemma}[The first restriction on the even letters]
At each even letter one can impose a single kernel condition \(z\in\ker\pi_i\) for which
\[
 \operatorname{im}\pi_i\text{ is generated by at most }m\text{ elements},\qquad z\in\ker\pi_i\ \Longrightarrow\ u_i(z)\in N_i(T_2),
\]
where \(u_i(z)\) is the corrected vector constructed in the proof. If \(m=0\) the assertion holds on
the full free family.
\end{lemma}
\begin{proof}
For \(x\in q^{-a}N_i(T_2)\), let
\(\partial x=X^{-1}(\gamma-1)x\) in the \'etale module.
It preserves \(q^{-a}N_i(T_2)\) and \(N_i(T_2)\).
Indeed the action on the invariant Wach lattice modulo \(X\)
is trivial, and \(\gamma(q)/q\) is a unit congruent to one
modulo \(X\). At a primitive root, put
\[
             c_j=\frac{1-\chi(\gamma)^{-j}}{\ell}.
\]
On the associated Laurent coefficient of degree \(-j\) the
divided operator acts by \(c_j\) since \(X=-\ell\) there.
For a Hodge jump \(u\), the quotient of the allowed pole lattice
by the unshifted lattice has degrees \(-u-k\), \(1\le k\le a\).

To see why one factor per distinct shift suffices, choose a basis
of the crystalline space adapted to the filtration. In each
direction the Laurent action is triangular, with distinct
diagonal entries indexed by the degrees in that direction.
The action of \(1/X(t)\) is scalar and the cyclotomic action
in the crystalline basis is substitution
\[
 t\mapsto\chi(\gamma)t.
\]

Thus the product vanishes separately on each direction even
when a shift repeats in another direction. This is a computation
over the coefficient field. The Hodge criterion at the cyclotomic
roots and saturation for \(\ell\) in \(q^{-a}N_i(T_2)/N_i(T_2)\)
give the same inclusion integrally.

Let \(J_i\) be the nonzero shifts that occur. They belong to
\([1-B_i,a-b_i]\cup[1-\eta,a+\eta]\).
By the valuation of powers of a principal unit,
\[
 v_\ell(c_j)=v_\ell(j)
\]
for \(j\ne0\). In the first
interval there are at most \((a+S)/\ell^k+1\) multiples of
\(\ell^k\), and no such nonzero multiples for
\(k>\log_\ell(2H)\). For the central interval split at zero
and apply the factorial bound on both sides.

Taking the maximum
over even letters, choose \(s_1\) and polynomials \(R_i(U)\)
proportional over \(O\) to \(\prod_{j\in J_i}(U-c_j)\) with
\[
 R_i(0)=\ell^{s_1},\qquad
 s_1\le(2a+2\eta+S)/(\ell-1)+L_0,\qquad
 \partial R_i(\partial)q^{-a}N_i(T_2)\subset N_i(T_2).
\]
Including \(\partial\) kills degree zero if that shift occurs.
The scaling of each polynomial is integral by the choice of the
maximum. Write \(R_i(U)=\ell^{s_1}+U R_i^*(U)\), and set
\[
 m_i=R_i(\partial)(q^{-a}\nu_i)+R_i^*(\partial)(z_i/X),
 \qquad u_i=\ell^R m_i-\ell^{s_1}\mu_i .
\]
Put \(v_i=\ell^R q^{-a}\nu_i-\mu_i\) and
\(d_i=\ell^R z_i/X+\partial\mu_i\).

By \eqref{eq:thick-extension-data},
\(q^a v_i,d_i\in N_i(T_2)\). The polynomial identity gives
\(u_i=R_i(\partial)v_i+R_i^*(\partial)d_i\). Thus
\(q^a u_i\in N_i(T_2)\) and \((\gamma-1)u_i\in XN_i(T_2)\).
Also
\[
 \ell^R\bigl(\ell^{s_1}z_i+(\gamma-1)m_i\bigr)
       =X\ell^{s_1}d_i+(\gamma-1)u_i .
\]
The bracket belongs to \(XN_i(T_2)\) by cancellation
of \(\ell^R\) in the integral \'etale module, where the Wach
lattice and its \(X\)-multiple are \(\ell\)-saturated. In coordinates this follows from the
description of the \'etale module as the completed localization
of the lattice, with injective reduction of the lattice modulo
the coefficient uniformizer.

At a primitive root, the only Laurent coefficient of \(u_i\)
that can obstruct membership in the Wach lattice is degree zero.
Writing \([t^j]\) for coefficient extraction at the root gives
\[
 [t^j](\gamma-1)u_i=(\chi(\gamma)^j-1)[t^j]u_i\qquad(j\ne0).
\]
The scalar on the right is nonzero. The degree-zero obstruction is therefore
the class in \(\mathrm{Fil}^{-a}D_i/\mathrm{Fil}^{0}D_i\).
All jumps in \([-a,0)\) are central since \(a<b_i\).
The quotient has dimension at most \(m\). Each obstruction
map is linear in the free family and its image is an
\(O\)-module generated by at most \(m\) elements.
Successively take these kernels for even \(i\). This leaves
a family on which every \(u_i\) belongs to \(N_i(T_2)\), using
one rank restriction per even letter. When \(m=0\) the entire
family survives.
\end{proof}

\paragraph{Cyclotomic correction at fixed precision.}
\label{par:thick-integral-correction}
Put
\[
 M=1+\lfloor h/(\ell-1)\rfloor,\qquad
 \sigma_1=\left\lceil 1+(M+\eta)/(\ell-1)\right\rceil .
\]
For a vector \(v\in XN_i(T)\) there is an integral correction
\(\tau(v)\in P_i^{(1)}(T)\) well-defined modulo
\(X^Mq^\eta P_i^{(1)}(T)\), with
\[
  \ell^{\sigma_1}\Phi v-(\gamma-1)\tau(v)
             \in X^Mq^\eta P_i^{(1)}(T).
\]
We use it at an even \(i\), where the odd destination gives
\(q^\eta P_i^{(1)}(T)\subset N_{i+1}(T)\).
There is also a two-step correction with precision \(Xq^{2H}\)
on \(P_i^{(2)}(T)\) and exponent
\(\sigma_*=\lceil1+2H/(\ell-1)^2\rceil\).

Here are the divisions that give these two claims. In either
case solve \((\gamma-1)d=v\) modulo a large power \(X^K\),
over the coefficient field. The degree-by-degree solution
of Paragraph~\ref{par:wach-cut-correction} puts the term in
degree \(y\) in \(\ell^{-e_y}X^yN_i(T)\), where
\[
 e_y=\lceil\ell y/(\ell-1)\rceil.
\]
Take \(d\) with zero
constant coefficient. Write \(J\) for the exponent of \(q\)
in the precision ideal and \(a_y=(\ell-1)y\). In one step
\(\phi(X)^y\) is a unit times \(X^yq^{a_y}\).

If \(a_y<J=\eta\), then \(e_y\le\sigma_1\), since
\(J\le (\ell-1)M\). If \(a_y\ge J\) and \(y<M\),
expand \(q^{a_y-J}=(X+\ell)^{a_y-J}\).
The terms which survive modulo \(X^Mq^J\) are divisible by
\(\ell^{\max(0,\ell y-J-M+1)}\). For \(e_y>\sigma_1\)
we have \(\ell y>J+M+\ell-1\); hence
\[
 \ell y-J-M+1\ge e_y-\sigma_1.
\]

For \(a_y\ge J\) and \(y\ge M\) the entire term vanishes.

For two steps use \(M=1\), \(J=2H\), and
\(\phi^2(X)^y\) associated to
\(X^yq^{a_y}p_1^{a_y}\).
If \(a_y\ge J\) the terms with \(y>0\) vanish. Otherwise
\(p_1\in(\ell,q^\ell)\), by reduction modulo \(\ell\).
After factoring out \(X^yq^{a_y}\), the terms of the binomial
expansion that can survive have \(q^\ell\) used at most
\(\lfloor(J-a_y-1)/\ell\rfloor\) times. If
\(e_y>\sigma_*\), then
\(a_y>J/\ell+(\ell-1)^2/\ell\), and
\[
 a_y-\frac{J-a_y-1}{\ell}
      \ge \frac{\ell a_y-J}{(\ell-1)^2}
      \ge e_y-\sigma_* .
\]
The last inequality uses the defining roundings. Thus the
required division is integral. Here one may take \(K>\max(M,J)\), since
\(\phi^j(X)^K\) then belongs to the precision ideal in both
cases.

The integral coefficient quotient by \(X^Mq^J\) is free
over \(O\) by monic Weierstrass division. For either free Frobenius span \(P\) and its precision submodule \(P_{\rm prec}\), saturation gives
\[
 (P_{\rm prec}[1/\ell])\cap P=P_{\rm prec}.
\]

The resulting representative of \(\ell^{\sigma_1}\Phi d\)
or \(\ell^{\sigma_*}\Phi^2d\) satisfies the claimed
cyclotomic congruence, by commutation and saturation.
Uniqueness of \(d\) modulo \(X^K\) over the coefficient field
follows from the nonzero diagonal \(\chi(\gamma)^{(\ell-1)y}-1\)
in positive degrees. Hence the corrections are linear modulo
the precision and natural in maps of crystalline lattices.
They may be made linear on a free family by choosing
representatives on its basis. In one step the correction to a
vector of \(X^MN_i(T)\) can be taken zero. The calculation of
\(d\) in that case starts at degree \(M\), and
\[
 (\ell-1)M\ge\eta.
\]
For \(v=(\gamma-1)w\) with
\(w\in XN_i(T)\) one can use \(d=w\).

\paragraph{Second restriction and the alternating sum.}
Write \(N_i=N_i(T_2)\) for this construction. For even \(c\)
set
\[
 \xi_c=\ell^{s_1}z_c+(\gamma-1)m_c\in XN_c.
\]

Choose \(t_c=\tau(\xi_c)\) linearly on the restricted family.
At the odd letter \(c+1\) put
\[
 L_c=\ell^{\sigma_1}\Phi m_c+
             \ell^{s_1+\sigma_1}y_{c+1}-t_c,\qquad
 \Xi_c=\ell^{s_1+\sigma_1}z_{c+1}+(\gamma-1)L_c .
\]
The Herr relation gives
\(\Xi_c=\ell^{\sigma_1}\Phi\xi_c-(\gamma-1)t_c\).

Consequently \(\Xi_c\in X^MN_{c+1}\). It lies more precisely
in \(X^Mq^\eta P_c^{(1)}(T_2)\).
To make the next difference integral, write
\[
 U_{c+1}=\ell^R L_c-\ell^{s_1+\sigma_1}\mu_{c+1}
        =\ell^{\sigma_1}\Phi u_c+\ell^{s_1+\sigma_1}A_{c+1}
           -\ell^R t_c .
\]
The signed height and \eqref{eq:thick-extension-data} show that \(q^\eta U_{c+1}\in N_{c+1}\).
Applying \(\gamma-1\) gives \(\ell^R\Xi_c\) minus
\(\ell^{s_1+\sigma_1}(\ell^R z_{c+1}+(\gamma-1)\mu_{c+1})\).
This is in \(N_{c+1}\). Its only possible Laurent obstruction
is again degree zero, this time in
\(\mathrm{Fil}^{-\eta}D_{c+1}/\mathrm{Fil}^0D_{c+1}\).
Only central jumps contribute, so taking one kernel of rank
at most \(m\) per odd destination makes each \(U_{c+1}\) integral.
There are at most \(n\) restrictions including the first ones.

On the resulting family define at each even \(c\)
\[
 w_c=\Phi L_{c-2}+\ell^{s_1+\sigma_1}y_c-\ell^{\sigma_1}m_c,
 \qquad Y_c=\Phi w_c+t_c .
\]
All indices are cyclic, including when the word has length two.
There are three useful tests on \(w_c\). First,
\[
 \ell^R w_c=\Phi U_{c-1}-\ell^{\sigma_1}u_c
                        +\ell^{s_1+\sigma_1}A_c .
\]
The height bounds place the right side, after multiplication
by \(q^{B_c}\), in \(N_c\). Multiplication instead by \(q^\eta\)
places it in \(P_{c-1}^{(1)}(T_2)\) after a possible
additional factor \(\ell^R\), by \eqref{eq:thick-extension-data}.
Both lattices are \(\ell\)-saturated inside the integral
\'etale module.

For the Frobenius span, reduce in a Wach basis
and extend scalars by the finite free Frobenius coefficient map.
The original basis stays independent over the residue Laurent
field because the linearized Frobenius on the \'etale module
is invertible. Cancel these powers of \(\ell\).
Second, \((\gamma-1)w_c=\Phi\Xi_{c-2}-\ell^{\sigma_1}\xi_c\)
belongs to \(XN_c\), since
\(\Phi(X^MN)\subset q X^MN\). Third,
\[
                (\gamma-1)Y_c=\Phi^2\Xi_{c-2}-\Xi_c
                                 \in X^MN_{c+1}.
\]
These calculations give exactly the pole, first pullback
and cyclotomic tests on \(w_c\) used by the even-letter estimates.

Define \(Z_c=(\Phi-1)^{-1}(\gamma-1)Y_c\) in \(X^MN\),
where the inverse exists by contraction on the product of all
letters. Then the sum of \((Y_c,Z_c)\) over even \(c\)
represents \(\ell^{s_1+\sigma_1}(y,z)\). Indeed let
\[
  B=\sum_{c\text{ even}}\bigl(L_c+\Phi L_c
                                  +\ell^{s_1+\sigma_1}y_{c+2}\bigr).
\]
Summing first components gives
\(\sum_cY_c=\ell^{s_1+\sigma_1}y+(\Phi-1)B\)
by invariance of the even labels under the shift by two.
For the second components
\(\ell^{s_1+\sigma_1}z+(\gamma-1)B
 =\sum_c(\Xi_c+\Phi\Xi_c)\in X^MN\).
The Herr relation and uniqueness of the inverse prove the same
coboundary identity for the whole pair.
\label{cor:inert-thick-herr-1}

\paragraph{Integral comparison with the lifted vectors.}
For an admissible \(w\) at an even \(c\), use the first
correction to form
\[
             V(w)=\ell^{\sigma_1}\Phi w
                          -\tau((\gamma-1)w).
\]
Its cyclotomic defect lies in \(X^MN\). The second component is uniquely determined at this precision by
\[
 z=(\Phi-1)^{-1}(\gamma-1)V(w),\qquad z\in X^MN.
\]

Its first component has denominators in finitely many
cyclotomic factors, by the height bounds, and the second
belongs to \(X^MN\). The crystalline criterion proved in
Section~\ref{sec:inert-long-orbit} therefore makes its
integral Herr class finite. A change of representative of
the correction changes both components by a coboundary.
Indeed if the first component of a pair is in \(X^MN\)
along with its second component, the first is the
\((\Phi-1)\)-image of a vector in \(X^MN\). Commutation
and injectivity on \(X^MN\) give the other component of
the coboundary. This also shows that the resulting class
depends linearly on \(w\) and is natural in the marked map.
Writing \([V(w)]\) for the resulting class, the last property of the correction gives
\[
 w\in XN_c\quad\Longrightarrow\quad[V(w)]=0.
\]
Thus integral congruence modulo \(XN_c\) is
sufficient for comparing classes.

For \(w=w_c\), the defect at letter \(c\) is
\(\Phi\Xi_{c-2}-\ell^{\sigma_1}\xi_c\).
Its first summand lies in \(X^MN_c\), on which the
correction is zero. Its second summand has correction
\(-\ell^{\sigma_1}t_c\) modulo the precision.
Thus
\[
 V(w_c)\equiv\ell^{\sigma_1}Y_c\pmod{X^MN}.
\]

The resulting finite classes, denoted by \(\kappa_c\), sum
to \(\ell^{s_1+2\sigma_1}\) times the original class.
The displayed cocycle uses \(\Phi-1\) and \(\gamma-1\)
in this order and has the displayed plus sign. If the
opposite convention for the extension class is used, replace
the sum of the source lifts by its negative.

In particular, if
\[
 \ell^e w_c-f(u_c')\in XN_c(T_2)
\]
for an admissible vector \(u_c'\) in the source, the
finite class attached to \(u_c'\) maps to
\(\ell^e\kappa_c\). All these comparisons take place on
the original rank block and the same map; they do not
require exactness of the Wach functor on the right or a
choice of integral splitting of the surjection.
\label{cor:inert-thick-herr-2}

\paragraph{Filtered comparison at a positive letter.}
\label{par:thick-filtered-comparison}
Write \(i\) for an odd letter, \(b=b_i\) and \(A=B_{i-1}\).
Suppose \(y\in N_{i-2}(T_2)\) has
\(z=\Phi^2\operatorname{sp}_{i-2}(y)\in\mathrm{Fil}^bD_i\).
\begin{proposition}[Integral preimage of the positive filtered class]
In the two-step height and jump range at \(i-2,i-1,i\), put
\[
 E'=\left\lfloor\frac{3a+S+2\eta-1}{\ell-1}
                         +\log_\ell H\right\rfloor .
\]
There exists a vector of the source Wach lattice at \(i\) whose
specialization maps to \(\ell^{E'+\sigma_*+A-b}z\).
The exponent \(E'\le E_G\) is independent of \(y\).
\end{proposition}
\begin{proof}

Use the two-step correction
of Paragraph~\ref{par:thick-integral-correction} to form
\[
 \lambda=\ell^{\sigma_*}\Phi^2 y-
                      \tau((\gamma-1)y)\in P_{i-2}^{(2)}(T_2).
\]
It has cyclotomic defect in \(Xq^{2H}P_{i-2}^{(2)}(T_2)\).
At a primitive root the span is regular in \(t\) and the
constant of \(\lambda\) is \(\ell^{\sigma_*}z\).
The correction has zero constant there because each transported
positive \(X\)-power vanishes at the root and its division modulo
\(Xq^{2H}\) preserves this. By the cyclotomic defect all
positive coefficients below \(2H\) vanish. The correction
likewise has zero specialization at \(X=0\).
These conclusions hold for the same integral representative
at every primitive root.

The signed height at \(i-1\) puts \(p_1^A\lambda\) in
\(P_{i-1}^{(1)}(T_2)\). Set \(x=q^{-b}p_1^A\lambda\).
Then
\[
 q^{b+\eta}x\in N_i(T_2)
\]
by the height at \(i\).
At a primitive root, the coefficients of \(x\) of degree
\(k<-b\) vanish. In degrees \(-b\le k<2H-b\) they lie
in the span of \(z\), hence in \(\mathrm{Fil}^{-k}D_i\).
The filtration is full for \(k\ge 2H-b\ge\eta\), since
all jumps at the odd letter are at least \(-\eta\).
The Hodge lattice test therefore gives \(x\in N_i(T_2)\).
The same signed height at \(i\) applied to the defect
gives
\[
           p_1^A(\gamma-1)\lambda\in Xq^{a+b}N_i(T_2),
\]
because \(2H\ge a+b+\eta\).
In particular
\(\operatorname{sp}_i(x)=\ell^{\sigma_*+A-b}z\).

It remains to lift a controlled multiple of \(x\) in the
integral Wach lattice. Let \(V_q=\gamma(q)/q\),
\(V_1=\gamma(p_1)/p_1=\phi(V_q)\).
Then \(V_q\) is one modulo \(X\) and \(\chi(\gamma)\) modulo \(q\).
The last equality follows by expanding at a primitive root
in the parameter \(t\) of degree one. Also
\(V_1\equiv1\pmod{Xq}\) by the formula for \(\phi(X)\).
The preceding defect equation implies
\[
 \gamma x=V_q^{-b}V_1^A x
                    \pmod{Xq^aN_i(T_2)}.
\]
Thus \(\partial x=d x\bmod q^aN_i(T_2)\) with
\(d=(V_q^{-b}V_1^A-1)/X\) integral and \(d\equiv c_b\bmod q\).

For \(g\in O[[X]]\), \(\gamma(g)\equiv g\bmod q\), and
the divided operator preserves \(q^aN_i(T_2)\).
Applying it to \(q^k g x\) shows that
\[
    (\partial-c_{b-k})q^k g x
           \in q^{k+1}O[[X]]x+q^aN_i(T_2)\qquad(0\le k<a).
\]
Indeed the multiplier on the surviving degree is
\(\chi(\gamma)^k c_b+(\chi(\gamma)^k-1)/(-\ell)=c_{b-k}\).
It follows that \(P_-(\partial)x\in q^aN_i(T_2)\), where
\(P_-(U)=\prod_{j=b-a+1}^{b}(U-c_j)\).
The image input of Paragraph~\ref{par:thick-initial-classes}
gives an integral source preimage of this vector.

For a second filter, first lift \(q^a x\) integrally and divide
the lift by \(q^a\) in the source \'etale module.
The resulting vector \(v\in q^{-a}N_i(T_1)\) maps to \(x\).
The shifted-jump calculation of the first restriction, now
on the odd source, shows that \(P_+(\partial)v\in N_i(T_1)\),
where
\[
 P_+(U)=\prod_{j\in[b+1,b+S+a]\,\cup\,[1-\eta,\eta+a]}(U-c_j).
\]
The two sets of indices for \(P_-\) and \(P_+\) are disjoint
because \(b\ge 2a+\eta\). This integer inequality follows
from \(b\ge 2A_0^{\rm real}+\eta+1\) and
\(a<A_0^{\rm real}+1\), with equality of \(a\) and the
real bound allowed when the latter is integral.

For a root index \(j\) of \(P_-\), the valuation of
\(P_+(c_j)P_-'(c_j)\) is the sum of \(v_\ell(j-j')\)
over the other indices. Indeed, if \(I_-\) and \(I_+\) denote the two index sets, factoring
the polynomials at the root gives
\[
 v_\ell(P_+(c_j)P_-'(c_j))
 =\sum_{j'\in(I_-\cup I_+)\setminus\{j\}} v_\ell(c_j-c_{j'})
 =\sum_{j'\in(I_-\cup I_+)\setminus\{j\}} v_\ell(j-j') .
\]
The indices in or above the first interval give factorials
with total weight \(2a+S-1\). In the central interval the
differences form an interval of length \(a+2\eta\) with
largest absolute value at most \(H\).
Thus the valuation is at most
\[
           \frac{3a+S+2\eta-1}{\ell-1}+\log_\ell H .
\]
Let \(E'\) be the floor of the displayed bound.

For the index set \(J_-\) of \(P_-\) take
\[
 B(U)=\sum_{j\in J_-}
 \frac{\ell^{E'}}{P_+(c_j)P_-'(c_j)}\,\frac{P_-(U)}{U-c_j}.
\]
The valuation bound gives \(B\in O[U]\). Since \(P_-\) is monic
and the numerator vanishes at every \(c_j\), there is \(C\in O[U]\)
with \(BP_++CP_-=\ell^{E'}\). Choose integral source preimages
\(v_+,v_-\) of \(P_+(\partial)x,P_-(\partial)x\), respectively,
as above. Stability under \(\partial\) then gives an integral vector
\(B(\partial)v_++C(\partial)v_-\). By commutation with \(f\), its
image is \(\ell^{E'}x\) and its specialization maps to
\(\ell^{E'+\sigma_*+A-b}z\).

The exponent \(E'\le E_G\) is independent of \(y\)
and of any choice of an integral splitting. This is the
filtered estimate used at each odd letter in the two
even-letter comparisons.
\label{cor:inert-thick-filtered-1}
\end{proof}


\paragraph{The zeroth coefficient at an even letter.}\label{par:thick-constant-data}
Use \(N_i\) for the target Wach lattice at \(i\), in the tame-invariant
coordinate ring. Put \(B_i=b_i+S\) at each letter.
Use the associate \(q=X+\ell\) of the first cyclotomic factor
and \(t\) for the logarithmic parameter at a primitive \(\ell\)-th root.
The parameter \(X\) of the
tame-invariant Wach ring is a unit in the completed de Rham ring at
that root. Write \(D_i\) for the \(i\)-th crystalline space and
\(\operatorname{sp}_i:N_i\to D_i\) for the reduction at \(X=0\).

Let \(\Phi\) also denote the induced Frobenius \(D_i\to D_{i+1}\).
The one-step Wach height at \(c-2\) gives
\[
\begin{aligned}
 P_{c-2}&=XN_{c-2}+q^{B_{c-2}}\langle\Phi N_{c-3}\rangle
                 \ \subset N_{c-2},\\
 \operatorname{sp}_{c-2}(P_{c-2})
       &=\ell^{B_{c-2}}\Phi\,\operatorname{sp}_{c-3}(N_{c-3}),
\end{aligned}
\]
Here the bracket denotes the span over the tame-invariant ring.
The equality follows by reduction at \(X=0\), since \(q(0)=\ell\)
and all scalar constants in \(O\) occur as reductions.
Take
\[
  C=\mathrm{Fil}^{b_{c-1}}D_{c-1},\qquad
  L_\eta=\mathrm{Fil}^{-\eta}D_c,\qquad L_0^{\rm fil}=\mathrm{Fil}^0 D_c .
\]
Both \(D_{c-1}/C\) and \(L_\eta/L_0^{\rm fil}\) have dimension at most
\(m\). The first quotient keeps only jumps in the central interval at
the odd letter. The second keeps only the negative central jumps at
the even letter.

We give the coefficient comparison with its quantitative inputs.
On a submodule \(A'\) of the free family of cocycles, suppose that each
\(w(a)\) is a tame-invariant vector at \(c\), linear in \(a\), with
\((\gamma-1)w(a)\in XN_c\) and
\[
 q^{\eta+B_c}w(a)\in N_c.
\]

Its normalized Laurent coefficients \(d_j(a)\) are obtained by
expanding at one primitive \(\ell\)-th root after dividing by the
expansion of \(q^{\eta+B_c}\), which is \(t^{\eta+B_c}\) times a
unit. They are linear in \(a\). The action of the principal generator
fixes the root and on the \(j\)-th logarithmic coefficient acts by
\(\chi(\gamma)^j\). For \(j\ne 0\), the coefficient of
\((\gamma-1)w(a)\) is in \(\mathrm{Fil}^{-j}D_c\) and
\(\chi(\gamma)^j-1\ne0\). Hence
\[
 d_j(a)\in\mathrm{Fil}^{-j}D_c
\]
for \(j\ne0\).

\paragraph{Transport to the first cyclotomic level.}
Use either Wach module and choose an integral vector \(z_1\)
at \(c-1\). For a positive precision \(D\), put
\(J=\lceil D/(\ell-1)\rceil\) and
\[
 \sigma_D=\lceil 1+(D+1)/(\ell-1)\rceil.
\]

On a Wach basis at \(c-1\), take a polynomial vector \(v\)
of degree \(<J\) with constant coefficient equal to that of
\(z_1\) and with \((\gamma-1)v\) in \(X^J\) times the Wach
module after inverting \(\ell\).
The degree-by-degree construction of
Paragraph~\ref{par:wach-cut-correction} bounds the denominator
in degree \(a\) by \(\ell^{A_a}\), \(A_a=a+v_\ell(a!)\).

Consider \(\ell^{\sigma_D}\Phi v\)
in the Frobenius span at \(c\) modulo \(Xq^D\).
In degree \(a\ge1\) its scalar before the coefficient denominators
is a unit times \(\ell^{\sigma_D}X^a q^{(\ell-1)a}\).
In the expansion \(X^a=X(q-\ell)^{a-1}\), the terms retained
modulo \(Xq^D\) have \(\ell\)-content at least
\(\max(0,\ell a-D)\). Since \(A_a\le\ell a/(\ell-1)\),
\[
 \sigma_D+\max(0,\ell a-D)\ \ge\ A_a .
\]
For \(\ell a\le D\) this follows from
\(\sigma_D\ge D/(\ell-1)\). For \(\ell a\ge D\) it
follows because \(\ell a/(\ell-1)-\ell a+D\) decreases
as \(a\) increases. Thus the transported vector has an
integral representative \(y\) in the Frobenius span modulo
\(Xq^D\) after inverting \(\ell\).

Since \(\phi(X)^J\) is divisible by \(Xq^D\), the representative
has \((\gamma-1)y\) in \(Xq^D\) times this span after inverting
\(\ell\). For the free Frobenius span \(P\), the quotient by the distinguished factor \(Xq^D\) is \(O\)-torsion-free. Consequently
\[
 (Xq^D P[1/\ell])\cap P=Xq^D P.
\]
This returns the rational divisibility to the integral span.
The construction is in the tame-invariant span. The coordinates
of \(y\) at \(X=0\) on the Frobenius basis are
\(\ell^{\sigma_D}\) times the coordinates of \(z_1\) at \(X=0\)
on the preceding basis.

At each primitive \(\ell\)-th root the
Frobenius span evaluates in the crystalline space with
coefficients in \(K_O[[t]]\). Its constant term is obtained
by \(\Phi\) from specialization at the preceding letter.
The \(\gamma-1\) congruence gives zero coefficients for
\(0<j<D\), because
\(\chi(\gamma)^j-1\ne0\) at each such degree.
Thus through degree \(D-1\) the expansion consists of the one
constant \(\ell^{\sigma_D}\Phi\operatorname{sp}_{c-1}(z_1)\).

The filtered crystalline comparison at the intervening odd letter
provides a number \(E'\le E_G\), where
\[
\sigma_2=1+\frac{2H_{\rm loc}}{(\ell-1)^2},\qquad
E_G=\frac{3A_0+S+2\eta}{\ell-1}+L_0 .
\]
For \(z\in \Phi^2\operatorname{sp}_{c-3}(N_{c-3})\cap C\) there
is \(z_1\in N_{1,c-1}\) whose specialization maps to \(\ell^k z\),
where \(k=\lceil E'+\sigma_2\rceil+B_{c-2}-b_{c-1}\).
Here \(N_{1,i}\) denotes the source Wach lattice.
Apply the transport in the source with \(D=B_c\).
The signed height gives \(q^{B_c}y\in N_{1,c}\) and
\((\gamma-1)y\in XN_{1,c}\). Its image has constant coefficient
\(\ell^\kappa\Phi z\) and all nonzero coefficients below \(B_c\)
vanish, with
\[
 \kappa\le (B_c+1)/(\ell-1)+3+E_G+\sigma_2+
                   B_{c-2}-b_{c-1},\qquad
 \kappa=\sigma_{B_c}+k.
\]

\paragraph{The lattice at an even letter.}
\begin{lemma}[Covolume after retaining the central directions]
\label{lem:thick-even-covolume}
At an even letter \(i\) of this word, start with
\(P_i=XN_i+q^{B_i}\langle\Phi N_{i-1}\rangle\).
Adjoining at most \(m\) vectors to \(P_i\) leaves normalized
colength at most \(C_D=rS+2mc_D\) in \(N_i\).
In particular
\[
 m=0\quad\Longrightarrow\quad\ell^S N_i\subseteq P_i.
\]
\end{lemma}

\begin{proof}
At an even \(i\), choose a Wach basis \(e_j\) at \(i-1\) whose
Frobenius specializations give an \(O\)-basis of the specialized
Frobenius lattice adapted to the Hodge filtration at \(i\).
Such a basis exists because the intersections of a lattice
with the subspaces of a flag are saturated. Each successive
quotient is free over \(O\). Lift the resulting basis from
the specialization to the Wach module, where it is a basis
by Nakayama's lemma. Write \(-\alpha_j\) for the jump of the
corresponding crystalline vector. Divide the indices according to the signed jumps of the adapted basis:
\[
 J_0=\{j:|\alpha_j|\le\eta\},\qquad |J_0|\le m,\qquad b_i\le\alpha_j\le B_i\quad(j\notin J_0).
\]


Let \(D=B_i+\eta\). For each central direction \(j\), apply the
first-level transport with precision \(D\) to \(e_j\), obtaining
\(y_j\) in the Frobenius span. The vector \(g_j=q^{\alpha_j}y_j\)
belongs to \(N_i\). To see this, at each root its constant
direction starts in degree \(\alpha_j\) since \(q=t\) times
a unit, and its crystalline direction is in
\(\mathrm{Fil}^{-\alpha_j}\). Its error starts in degree
\[
 D+\alpha_j\ge B_i,
\]
where the filtration is full.
The \(q\)-power denominator is removed by the same integral
Weierstrass test as above.

For a Wach basis \((e_j)\) at \(i\), factoring out the common powers from the columns gives
\[
 \det(q^{B_i}\Phi e_1,\ldots,q^{B_i}\Phi e_r)
       =u q^{rB_i-\sum_j\alpha_j},\qquad u\in O[[X]]^\times .
\] The exponent follows at the
cyclotomic level because the Hodge lattice in the adapted
directions has orders \(\alpha_j\).
The two-sided signed height makes this matrix invertible
after inverting \(q\) over \(O[[X]]\), and \(q=X+\ell\)
generates a prime ideal. These facts give the asserted unit
and power over the integral ring. Leave the noncentral
columns as they are and replace each central one by \(g_j\).
Relative to the old columns the replacement is given by the
transport matrix with the central columns multiplied by
\(q^{\alpha_j-B_i}\). At \(X=0\), before the individual
multipliers the transport matrix has determinant
\(\ell^{m'\sigma_D}\),
where \(m'\le m\) is the number of central columns.

If \(C_{\rm old},C_{\rm new}\) denote the specialized
column matrices, then
\[
 \det C_{\rm new}=\det C_{\rm old}\cdot
       \ell^{m'\sigma_D+\sum_{j\ {\rm central}}(\alpha_j-B_i)}.
\]
Its normalized valuation is therefore
\[
\sum_{j\ {\rm noncentral}}(B_i-\alpha_j)+m'\sigma_D
 \le rS+2mc_D.
\]
Indeed \(D\le H_{\rm loc}\) gives
\(\sigma_D\le\lceil c_D\rceil\le2\lfloor c_D\rfloor\).
The noncentral columns belong to \(P_i\), and adding the
at most \(m\) central classes to \(P_i\) thus leaves a quotient
of normalized length at most \(C_D\).
When \(m=0\) the completed Hodge lattice is contained in
\(q^{b_i}\) times the completed Frobenius span. Weierstrass
saturation gives this inclusion integrally. Multiplication
by \(q^S\) then lands in \(P_i\), and
\(q^S\equiv\ell^S\bmod X\) gives the sharper
\(\ell^S N_i\subset P_i\).
\end{proof}

\paragraph{Recovering the target coefficient integrally.}
\begin{lemma}[An integral representative of the constant coefficient]
Use Lemma~\ref{hyp:alternating-thick-word} at an even letter \(c\).
Let \(w(a)\) be the tame-invariant vector with
\((\gamma-1)w(a)\in XN_c\) and \(q^{\eta+B_c}w(a)\in N_c\),
with the coefficient \(d_0(a)\) and filtered space \(L_\eta\)
of Paragraph~\ref{par:thick-constant-data}. There is
\(x\in N_{c-2}\) such that
\[
 \ell^{e_c}d_0(a)-\Phi^2\operatorname{sp}_{c-2}(x)\in L_\eta,\qquad
 e_c\le\sigma_2+B_{c-1}-b_c/\ell+L_0+3.
\]
\end{lemma}
\begin{proof}
Write \(\overline d\) for the class of \(d_0(a)\) modulo \(L_\eta\)
and \(y_0=B_{c-1}\). The coefficients of \(w(a)\)
in nonzero degrees \(j<b_c\) lie in \(L_\eta\). For \(j\le\eta\)
this follows from \(\mathrm{Fil}^{-j}\subset L_\eta\).
For \(\eta<j<b_c\) the two filtrations agree, by the gap
between the negative and central jumps at the even letter.
Take \(D=2b_c+S+\eta\), \(P=\phi(q)\), and
\(\sigma_H=\lceil 1+D/(\ell-1)^2\rceil\).
We first obtain the integral coefficient relation
\[
 \ell^{\sigma_H}q^\eta P^{y_0}\overline d
       \equiv \sum_j a_j(q)\mu_j \pmod{q^{b_c+\eta}} .
\]
The coefficients \(a_j\) are in \(O[[q]]\) and each \(\mu_j\) is
the image of \(\Phi^2\operatorname{sp}_{c-2}(x_j)\) for some
\(x_j\in N_{c-2}\).

The pole input \(q^\eta w(a)\in\langle\Phi N_{c-1}\rangle\)
and the opposite height at \(c-1\) give
\(q^\eta P^{y_0}w(a)\in\langle\Phi^2 N_{c-2}\rangle\).
Transport an entire truncated horizontal basis and its inverse
from \(c-2\), now by two Frobenius steps. We give the scalar
integrality for both matrices. Truncate at degree
\(J=\lceil D/(\ell-1)\rceil\) so the first omitted degree is at
least the degree needed for the \(D\)-jet test. The inverse has the
same denominator bound \(A_a=a+v_\ell(a!)\) in degree \(a\),
since
\[
 A_{i+j}-A_i-A_j=v_\ell\binom{i+j}{i}\ge0.
\]
Iterating this inequality controls every product in the inverse matrix series.

The degree-\(a\) scalar \((\phi^2 X)^a\) is a unit times
\[
        X^a q^{(\ell-1)a}P^{(\ell-1)a}.
\]
In the \(q\)-parameter \(P=\ell\epsilon+q^\ell v\) with
\(\epsilon,v\in O[[q]]^\times\), by
\(\phi(X)=Xq^{\ell-1}\) times a unit and \(X=q-\ell\).
In the binomial expansion of \(P^{(\ell-1)a}\), an index \(j\)
contributes a factor \(q^{\ell j}\ell^{(\ell-1)a-j}\).
For \(a\ge1\), the term with index \(j\) is divisible by \(Xq^D\) as soon as \((\ell-1)a+\ell j\ge D\). Thus a surviving term has
\[
 j<\frac{D-(\ell-1)a}{\ell},\qquad
 (\ell-1)a-j>(\ell-1)a(1+1/\ell)-D/\ell .
\]
The constant term in the transported matrix has no denominator. The retained terms
have \(\ell\)-content at least
\[
\max\bigl(0,(\ell-1)a(1+1/\ell)-D/\ell\bigr).
\]

Put \(a_*=D/((\ell-1)(\ell+1))\). For \(a\le a_*\) the cost bound
\(\ell a/(\ell-1)\) increases. For \(a\ge a_*\) one subtracts the
linear content term, so the slope becomes
\[
 \frac{\ell}{\ell-1}-(\ell-1)(1+1/\ell)<0\qquad(\ell\ge3).
\]
Extending the cost bound from integral degrees to real arguments, its maximum is attained at the junction of the two linear bounds:
\[
 \max_{a\in\mathbb R_{\ge0}}\left\{\frac{\ell a}{\ell-1}
       -\max\left(0,(\ell-1)a(1+1/\ell)-D/\ell\right)\right\}
 =\frac{\ell a_*}{\ell-1}
 =\frac{\ell D}{(\ell+1)(\ell-1)^2}
 \le \frac{D}{(\ell-1)^2}.
\]

Thus \(\ell^{\sigma_H}\) clears both transported matrices
modulo \(Xq^D\). Denote their integral representatives by \(Y,Y'\),
in the order of the horizontal basis and its inverse. The first
omitted degree contributes nothing because \(Xq^D\mid\phi^2(X)^J\).
For the integral coordinate vector \(\mathbf v\) of
\(q^\eta P^{y_0}w(a)\) in the two-step Frobenius span, put
\begin{equation}\label{eq:thick-two-step-inverse}
 \mathbf a=Y'\mathbf v,\qquad
 Y\mathbf a=YY'\mathbf v\equiv\ell^{2\sigma_H}\mathbf v
                                              \pmod{Xq^D}.
\end{equation}

Evaluate at a primitive \(\ell\)-th root. Two Frobenius steps take
the initial cyclotomic argument to the identity. The columns of
\(Y\) in the Frobenius basis are \(\gamma\)-invariant to precision
\(Xq^D\), as in the one-step construction. Their jets through
order \(D-1\) therefore equal the constants
\(\ell^{\sigma_H}\Phi^2\operatorname{sp}_{c-2}(x_j)\).
In the quotient by \(L_\eta\) these become
\(\ell^{\sigma_H}\mu_j\).

Write a bar for projection to this quotient.
The jump gap proved above gives
\[
 \overline{w(a)}\equiv\overline d\pmod{q^{b_c}},\qquad
 q^\eta P^{y_0}\overline{w(a)}
    \equiv q^\eta P^{y_0}\overline d\pmod{q^{b_c+\eta}} .
\]
These congruences use \((q)=(t)\) after evaluation.
Divide the evaluation of \eqref{eq:thick-two-step-inverse} by \(\ell^{\sigma_H}\) over
the coefficient field. Since \(D\ge b_c+\eta\), it gives the claimed
relation with the integral coordinates \(a_j(q)\) of \(\mathbf a\).
The translation \(q=X+\ell\) in the integral series is valid by
completeness of \(O\).

Let \(L\) be the \(O\)-span of the \(\mu_j\). It is finite free.
The coefficient of \(q^\eta\) shows
\[
 \overline d\in K_O L
\]
because \(P(0)\ne0\). Pass to an \(O\)-basis of \(L\), write
one coordinate of \(\overline d\) as \(\theta\), and divide the
resulting congruence by \(q^\eta\). This gives
\(\ell^{\sigma_H}\theta P^{y_0}\) integral modulo \(q^{b_c}\).
Put \(j_*=\lfloor(b_c-1)/\ell\rfloor\le y_0\), the inequality
following from \(b_{c-1}\ge h/2\). Among
\(j=0,\ldots,j_*\), choose the least index \(j_1\) minimizing
\(v_\ell\binom{y_0}{j}+y_0-j\).

Larger indices contribute only at degrees \(>\ell j_1\),
and smaller ones have strictly larger valuation by the minimal
choice of \(j_1\). The coefficient from \(j_1\) is a unit
multiple of \(\binom{y_0}{j_1}\ell^{y_0-j_1}\), since the
constant terms of the two units are units. Thus
\[
 \begin{aligned}
 v_\ell([q^{\ell j_1}]P^{y_0})
   &= y_0-j_1+v_\ell\binom{y_0}{j_1}\\
   &\le y_0-j_*+v_\ell\binom{y_0}{j_*}
    \le y_0-b_c/\ell+1+\log_\ell y_0 .
 \end{aligned}
\]
The last inequality uses the factorial floor sum and
\(j_*=\lfloor(b_c-1)/\ell\rfloor\).
It follows that every coordinate of \(\ell^{e_c}\overline d\)
is integral for
\[
e_c=\sigma_H+\lceil y_0-b_c/\ell+2+\log_\ell y_0\rceil .
\]
Lifting the resulting \(O\)-linear combination of the \(\mu_j\)
gives \(x\in N_{c-2}\) with
\[
 \ell^{e_c}d_0(a)-\Phi^2\operatorname{sp}_{c-2}(x)\in L_\eta .
\]
Since \(\sigma_H\le\sigma_2+1\) and \(y_0\le H_{\rm loc}\),
this exponent satisfies
\[
 e_c\le \sigma_2+B_{c-1}-b_c/\ell+L_0+3 .
\]

\end{proof}
\paragraph{Restrictions on the coefficients at the even letters.}
\begin{lemma}[Reduction to a Wach residual]
At an even letter let \(N_g,N_d\) count the restrictions by generator number and by length, respectively. To reduce to an integral Wach residual it suffices to take
\[
 N_g=3,\qquad N_d=1,
\]
with each of the three generator quotients generated by at most \(m\) elements. For \(m=0\) the reduction holds on the
whole family at exponent \(E_{\rm lev}+S\).
\end{lemma}
\begin{proof}
The height cost cancels before taking a bound over the letters.
Indeed \(e_c+\kappa-B_{c-2}\le B_{\rm lev}\), where
\[
 B_{\rm lev}=\frac{h}{\ell(\ell-1)}+\frac{S+1}{\ell-1}
              +S+E_G+2\sigma_2+6+L_0 .
\]
The calculation uses
\((B_c+1)/(\ell-1)-b_c/\ell
 =b_c/(\ell(\ell-1))+(S+1)/(\ell-1)\).
Put \(E_{\rm lev}=\lfloor B_{\rm lev}\rfloor\).
As \(e_c+\kappa-B_{c-2}\) is integral, the bound gives
\[
  E_{\rm lev}-e_c-\kappa+B_{c-2}\ \ge0,\qquad
  \ell^{E_{\rm lev}}d_0(a)
   \equiv \ell^{\kappa-B_{c-2}}\Phi^2\operatorname{sp}_{c-2}(x)
                      \pmod{L_\eta}
\]
for \(x\in N_{c-2}\). The correction is obtained by multiplying
the earlier choice of \(x\) by the nonnegative power
\(\ell^{E_{\rm lev}-e_c-\kappa+B_{c-2}}\).

Let \(K_\eta\) be the kernel of
\(\Phi^2\operatorname{sp}_{c-2}\bmod L_\eta\) on \(N_{c-2}\).
Multiplication by \(\ell^{\kappa-B_{c-2}}\) is invertible
on the coefficient-field quotient. Hence two choices of \(x\)
differ by \(K_\eta\), and addition of congruences gives an
\(O\)-linear map to \(N_{c-2}/K_\eta\). Composing with the
further quotient gives
\[
 A'\ \longrightarrow\ N_{c-2}/K_\eta\
       \longrightarrow\ N_{c-2}/(P_{c-2}+K_\eta).
\]
Write \(Q=N_{c-2}/(P_{c-2}+K_\eta)\) and \(I\subset Q\) for
the image of \(A'\). The kernel consists of the elements for which
one can choose \(x\in P_{c-2}\).

By Lemma~\ref{lem:thick-even-covolume}, there is a submodule
\(U\subset Q\) generated by at most \(m\) distinguished classes
with normalized length \(\ell_{\rm norm}(Q/U)\le C_D\).
The two restrictions on \(A'\) have successive quotients read from
\begin{equation}\label{eq:thick-restriction-quotients}
 0\longrightarrow I\cap U\longrightarrow I
       \longrightarrow (I+U)/U\longrightarrow0,\qquad
 \ell_{\rm norm}((I+U)/U)\le C_D .
\end{equation}
For the first term of \eqref{eq:thick-restriction-quotients}, the preimage of \(I\cap U\) in a free cover \(O^m\to U\) is free
of rank at most \(m\). Its image generates \(I\cap U\); killing that
image costs one generator restriction.

On the resulting kernel the new coefficients satisfy
\[
 \ell^{E_{\rm lev}}d_0(a)\equiv\ell^\kappa\Phi z\pmod{L_\eta},
 \qquad z\in Z:=\Phi^2\operatorname{sp}_{c-3}(N_{c-3}).
\]
Set \(K'=\ker(Z\xrightarrow{\,\Phi\,}D_c/L_\eta)\), so \(z\) is determined modulo
\(K'\). Requiring a representative in \(C\) costs at most \(m\)
generators: the quotient \(Z/((Z\cap C)+K')\) is an image of
\[
 Z/(Z\cap C)\hookrightarrow D_{c-1}/C,\qquad
 \dim_K(D_{c-1}/C)\le m ,
\]
where \(K\) denotes the coefficient field. The source of the
injection is finitely generated over \(O\), since \(Z\) comes from
a Wach lattice modulo \(X\). It is therefore free of rank at most
\(m\), and submodules of its quotient also need at most \(m\)
generators. For \(m=0\) the inclusion
\[
 \ell^S N_{c-2}\subset P_{c-2}
\]
allows \(\ell^S x\) as representative and
\(C=D_{c-1}\).

Choose the source vectors \(y\) from the first-level construction on
an \(O\)-basis of the restricted family, then extend linearly.
The constant of \(\ell^{E_{\rm lev}}w(a)-f(y(a))\) lies in
\(L_\eta\). Requiring it to lie in \(L_0^{\rm fil}\) costs one more
generator restriction of size at most \(m\). When \(m=0\) these
two filtration steps agree. For \(j\ne0\) the difference of
the coefficients lies in \(\mathrm{Fil}^{-j}D_c\), since the
source image has vanishing coefficients for \(j<B_c\) and
\[
 \mathrm{Fil}^{-j}D_c=D_c
\]
for \(j\ge B_c\).
Tame transport carries this property from the chosen root to
all primitive \(\ell\)-th roots. The rational cyclotomic
evaluation test now puts the vector difference in \(N_c\).

For the integral test, start with the known denominator
\(q^{\eta+B_c}\) and divide its coordinates in a Wach basis by the
distinguished polynomial associated to that power. At each root the
required coefficient vanishing follows from membership in the
completed Hodge lattice. The remainder is therefore zero by
Weierstrass division. This proves the reduction to a Wach vector
residual with three generator restrictions and one length restriction.
For \(m=0\) it works on the whole family with exponent
\(E_{\rm lev}+S\).

\end{proof}
\paragraph{The residual modulo \(X\).}
\begin{lemma}[Lifting the remaining class modulo \(X\)]
Three generator restrictions and two length restrictions lift the residual modulo \(XN_c\) after multiplication by \(\ell^{e_6}\), where
\[
 e_6=S+\lceil E_G+\sigma_2\rceil.
\] If \(m=0\), exponent
\(e_6+2S\) suffices on the entire family.
\end{lemma}
\begin{proof}
Use the same filtered crystalline comparison and
Lemma~\ref{lem:thick-even-covolume} at both even letters \(c,c-2\). Let \(v\) be the
integral Wach residual on the family.
Restrict first to \(v\in P_c\), costing one length and one
generator restriction, as above. There is a vector \(v_1\in N_{c-1}\)
with \(\operatorname{sp}_c(v)=\ell^{B_c}\Phi\operatorname{sp}_{c-1}(v_1)\).
At the odd letter the opposite Wach height gives
\(q^{B_{c-1}}N_{c-1}\subset\langle\Phi N_{c-2}\rangle\).
Specializing at \(X=0\) gives \(v_2\in N_{c-2}\) with
\[
 \ell^{B_{c-1}}\operatorname{sp}_{c-1}(v_1) =\Phi\operatorname{sp}_{c-2}(v_2).
\]

The specializations of \(v_1,v_2\) are linear in the family,
because Frobenius is invertible on the crystalline spaces and the
scaling factors are nonzero. The kernel of specialization is
\(XN_{c-2}\), so the map \(v_2\bmod P_{c-2}\) is well defined.
Restricting it to zero costs one length and one generator
restriction. Then for some \(v_3\in N_{c-3}\) one has
\(\operatorname{sp}_{c-2}(v_2)
=\ell^{B_{c-2}}\Phi\operatorname{sp}_{c-3}(v_3)\).
Restrict to the kernel of the image of
\(z=\Phi^2\operatorname{sp}_{c-3}(v_3)\) in \(D_{c-1}/C\).
This is one generator restriction of size at most \(m\).

Now \(e_6=S+\lceil E_G+\sigma_2\rceil\) satisfies
\[
 e_6+B_{c-2}=u_0+k+B_{c-1},\qquad
 u_0=\lceil E_G+\sigma_2\rceil-\lceil E'+\sigma_2\rceil\ge0 .
\]
Choose \(z_1\in N_{1,c-1}\) with
\(f(\operatorname{sp}_{c-1}z_1)=\ell^k z\).
Combining the three specialized equalities gives
\[
 \operatorname{sp}_c(\ell^{e_6}v)
   = f\bigl(\operatorname{sp}_c(\ell^{u_0}q^{B_c}\Phi z_1)\bigr).
\]
The vector \(\ell^{u_0}q^{B_c}\Phi z_1\) belongs to \(N_{1,c}\)
by the source height, so it is tame invariant and has
\(\gamma-1\) difference in \(XN_{1,c}\).
The residual is therefore lifted modulo \(XN_c\) at the cost of
three generator and two length restrictions. When \(m=0\), one
uses \(\ell^S N_i\subset P_i\) twice and \(C=D_{c-1}\).
This gives a lift on the entire family with exponent \(e_6+2S\).

\end{proof}
\begin{proposition}[Finite-class loss on an alternating word of larger jumps]\label{prop:thick-inert-colength}
For the marked crystalline surjection and signed jump word of
Lemma~\ref{hyp:alternating-thick-word}, the chain uses at most \(4n\) generator restrictions of
size \(m\) and \(3n/2\) length steps. Let \(A_{\rm end}\) be the last submodule in the restricted chain and
\(\pi:A\to M_2=H^1_f(F,T_2)\) the chosen surjection. There is an exponent \(e\) such that
\[
 \ell^e\pi(A_{\rm end})\subseteq f_*H^1_f(F,T_1),\qquad
 e\le \frac{32H_{\rm loc}}{\ell^2}+4S+\frac{10\eta}{\ell-1}+25L_0 .
\]
For \(m=0\) the construction gives a lift on the full finite
group with the additional allowance \(3S\). With the separately
available coarse lift, \eqref{eq:thick-chain-mass} converts these
losses to the normalized place colength.
\end{proposition}

\begin{proof}
Paragraphs~\ref{par:thick-initial-classes} and
\ref{par:thick-integral-correction} decompose the target finite
classes after \(s=s_1+2\sigma_1\), with at most \(n\) generator
restrictions of size \(m\). They realize the vectors \(w_c\)
on the target and compare their finite classes under the
same marked map modulo \(XN_c\).

\proofstep{Restrict the same free family at all even letters, then sum the cut classes.}
Now apply the letter
comparisons successively on the submodules obtained so far.
They are stable under restriction
of the family. The first lift \(y\) and the lift of the residual
combine by adding \(\ell^{e_6}y\) to the latter source vector. There
are \(n/2\) even letters. Counting at every other letter gives the respective maxima
\[
 n+6(n/2)=4n\quad\text{generator steps},\qquad 3(n/2)=3n/2\quad\text{length steps}.
\]
Each cut class has a finite lift after the same power.
Summing their classes lifts the scaled original class, since the
decomposition holds modulo coboundaries. This costs
\(e=s+E_{\rm lev}+e_6\), independently of the number of letters.
For \(m=0\) the extra powers on the two comparisons give \(e+3S\)
on the entire finite group.

\proofstep{Combine the decomposition and residual losses, using \(\ell\ge11\).}
The successive restrictions have been imposed on one family of lifts. Their losses can now be added without counting a different family at each even position.

Here is the numerical bound from this construction. Put
\(u=1/(\ell-1)\) and
\(\sigma=1+(\lfloor hu\rfloor+1+\eta)u\).
The resulting bounds for the decomposition and its remainder are
\[
\begin{aligned}
 s&\le(2A_0+2\eta+S)u+L_0+2\sigma+2,\\
 e_6&\le S+E_G+\sigma_2+L_0 .
\end{aligned}
\]
Adding these bounds and the bound for \(E_{\rm lev}\), using
\(\sigma\le1+hu^2+u+\eta u\), gives
\[
 e\le 8A_0u+8\eta u+(4S+3)u+9H_{\rm loc}u^2+2S+5L_0+13 .
\]
Since \(u\le1/10\) and
\(A_0u\le (2+u)H_{\rm loc}u^2+(4+L_0)u\), we obtain
\[
 e\le25.8H_{\rm loc}u^2+8\eta u+2.4S+5.8L_0+16.5
   \le \frac{32H_{\rm loc}}{\ell^2}+4S+
              \frac{10\eta}{\ell-1}+25L_0 ,
\]
using \(L_0\ge1\) and \(\ell/(\ell-1)\le 11/10\).
With the common coarse finite-class lift \(e'\),
\eqref{eq:thick-chain-mass} gives the bound for sequences with larger jumps.

The target coefficient bound uses the estimate on the rank of the kernel. The filtered crystalline estimate gives the realization at the first level. Apply the two lattice bounds for even letters in their respective cases.
 Combine the squared Frobenius bound on the residual lattice for an even letter with the coefficient lift.
 The resulting lattice bounds give the length at the chosen prime.
The local argument applies the integral Wach and Herr comparisons
to the marked crystalline models and uses the corresponding
coarse class lift for the generator restrictions.
\end{proof}

